% id: 127-04
% book: 베이직쎈 확률과 통계 · 07 통계적 추정
% section: 유형 077 표본평균의 평균, 분산, 표준편차; 모평균, 모표준편차가 주어진 경우
% figure: none
% answer: 
% answer_source: 
% variant_level: 0 (원본 전사)
% 이 파일은 build.mjs 가 items.json 에서 만든다. 고칠 때는 items.json 을 고친다.
\smash{\raisebox{1.5mm}{\dmkichul{베이직쎈 확률과 통계 127쪽 04번}}}
모평균이 70, 모표준편차가 5인 모집단에서 크기가 100인 표본을 임의추출할 때, 표본평균 $\overline{X}$에 대하여 옳은 것만을 \textbf{보기}에서 있는 대로 고른 것은? \bogi{ㄱ. $\overline{X}$는 확률변수이다.\\ ㄴ. $\overline{X}$의 평균은 70이다.\\ ㄷ. $\overline{X}$의 분산은 $\dfrac{1}{2}$이다.}
\dmchoicesauto{}{ㄱ}{ㄱ, ㄴ}{ㄱ, ㄷ}{ㄴ, ㄷ}{ㄱ, ㄴ, ㄷ}
